% ============================================================
% Brief research slides (4 pages): research question, literature,
% why it matters, empirical strategy. Companion to main.tex.
% Compile: pdflatex slides.tex
% ============================================================
\documentclass[aspectratio=169,11pt]{beamer}
\usetheme{default}
\usecolortheme{default}
\setbeamertemplate{navigation symbols}{}
\setbeamertemplate{footline}{}
\setbeamercolor{background canvas}{bg=white}
\setbeamerfont{frametitle}{series=\bfseries}
\setbeamercolor{frametitle}{fg=black}
\setbeamercolor{framesubtitle}{fg=black!55}

\usepackage[T1]{fontenc}
\usepackage{lmodern}
\usepackage{array}
\usepackage{amsmath}
\usepackage{xcolor}

\definecolor{accent}{HTML}{B5622A}
\definecolor{inksoft}{HTML}{4B554F}
\definecolor{protective}{HTML}{4F7A5B}
\definecolor{mixed}{HTML}{B3872F}
\definecolor{harmful}{HTML}{A83E32}
\definecolor{precedent}{HTML}{626B64}
\definecolor{panel}{HTML}{EEEEE7}

\newcommand{\littag}[2]{\textcolor{#1}{\scriptsize\bfseries #2}}

\title{}
\author{}
\date{}

\begin{document}

% ---------------- SLIDE 1: RESEARCH QUESTION ----------------
\begin{frame}
\frametitle{Research question}
\vspace{0.6cm}
{\Large\bfseries Does proximity to a legal recreational-marijuana market reduce drug-overdose mortality in states that haven't legalized it themselves?}

\vspace{1cm}
\begin{block}{}
\small We measure \textbf{proximity} as the distance from a never-legal county to the nearest county with an open retail market, tracked as that frontier moves outward with each new state's launch, \textbf{2020--2025}.
\end{block}

\vfill
{\scriptsize\color{inksoft} Cross-border spillover design \hfill Marijuana Legalization \& Overdose Mortality}
\end{frame}

% ---------------- SLIDE 2: LITERATURE ----------------
\begin{frame}
\frametitle{What does the literature say?}
\framesubtitle{Inconclusive --- sign, mechanism, and margin all disagree}

\footnotesize
\renewcommand{\arraystretch}{0.95}
\begin{tabular}{@{}p{3.0cm}p{7.6cm}r@{}}
\textbf{Bachhuber et al.} \newline {\scriptsize 2014 \textperiodcentered\ JAMA Internal Medicine} &
State medical-marijuana laws associated with \emph{fewer} opioid overdose deaths. &
\littag{protective}{Protective} \\
\textbf{Shover et al.} \newline {\scriptsize 2019 \textperiodcentered\ PNAS} &
Re-running Bachhuber's design on more years of data --- the association \emph{reverses sign}. &
\littag{mixed}{Reverses} \\
\textbf{Powell, Pacula \& Jacobson} \newline {\scriptsize 2018 \textperiodcentered\ J.\ Health Econ.} &
Any effect is specific to \emph{dispensary-based} designs --- legal status alone isn't enough. &
\littag{mixed}{Margin-specific} \\
\textbf{Mathur \& Ruhm} \newline {\scriptsize 2022 \textperiodcentered\ NBER WP 29802} &
Retail dispensaries linked to \emph{higher}, not lower, opioid mortality. &
\littag{harmful}{Opposite sign} \\
\textbf{Denkyirah et al.} \newline {\scriptsize 2026 \textperiodcentered\ Southern Econ.\ J.} &
State-level design on \emph{recreational} legalization: $\approx$3.5 fewer deaths per 100,000. &
\littag{protective}{Protective} \\
\textbf{Hao \& Cowan} \newline {\scriptsize 2020 \textperiodcentered\ Economic Inquiry} &
Cross-border spillover design precedent --- arrests, not mortality, as the outcome. &
\littag{precedent}{Design precedent} \\
\end{tabular}

\vspace{0.05cm}
\begin{block}{}
\scriptsize\color{black} No consensus on direction (protective vs.\ harmful), the relevant policy margin (medical vs.\ dispensary vs.\ recreational), or whether the effect is even stable over time.
\end{block}
\end{frame}

% ---------------- SLIDE 3: WHY IT MATTERS ----------------
\begin{frame}
\frametitle{Why does it matter}
\framesubtitle{Three reasons this question is worth answering carefully}

\vspace{0.3cm}
\begin{columns}[t,totalwidth=\textwidth]

\column{0.32\textwidth}
\begin{block}{\scriptsize\textcolor{accent}{AN OPEN DISPUTE} \\ \normalsize The literature genuinely disagrees}
\footnotesize Existing estimates disagree on sign, on which policy margin matters, and on whether the effect is even stable over time --- this isn't a settled question with one outlier.
\end{block}

\column{0.32\textwidth}
\begin{block}{\scriptsize\textcolor{accent}{A SUBSTITUTION QUESTION} \\ \normalsize Do people switch away from riskier drugs?}
\footnotesize If legal access lets high-risk individuals substitute marijuana for opioids, that's a direct, testable channel for harm-reduction policy.
\end{block}

\column{0.32\textwidth}
\begin{block}{\scriptsize\textcolor{accent}{A POLICY-SPILLOVER QUESTION} \\ \normalsize Effects reach beyond the legalizing state}
\footnotesize Because the outcome is measured in \emph{never-legal} counties, the design speaks directly to whether one state's policy affects health outcomes next door.
\end{block}

\end{columns}
\end{frame}

% ---------------- SLIDE 4: METHODOLOGY ----------------
\begin{frame}
\frametitle{How we answer this question}
\framesubtitle{A continuous distance-to-market design, on a stacked county panel}

\begin{columns}[t,totalwidth=\textwidth]

\column{0.48\textwidth}
{\small\bfseries Treatment \& sample}
\vspace{0.2cm}
\begin{itemize}
\small
\item[\textcolor{accent}{$\bullet$}] \textbf{Primary treatment (continuous):} miles from a county to the nearest county with an \emph{open retail market}, recomputed each year as more states launch.
\item[\textcolor{accent}{$\bullet$}] \textbf{Sample:} counties in the 31 states not yet legal as of the 2020 baseline --- a ``control reservoir'' whose distance to the frontier keeps shrinking.
\item[\textcolor{accent}{$\bullet$}] \textbf{Secondary measure:} a binary ``switcher'' indicator (distance crosses 50 miles), reported alongside, not in place of, the continuous measure.
\end{itemize}

\column{0.48\textwidth}
{\small\bfseries Specification}
\vspace{0.2cm}
\small Stacked cohort panel, one row per county per one-year transition it's still ``at risk'' in --- a county drops out once it switches, so later transitions never treat an already-treated county as a control.

\vspace{0.35cm}
\begin{block}{}
\small
$\Delta\text{Overdose}_{it} = \beta\cdot\Delta\text{Distance}_{it} + \text{Transition}_t + \text{Baseline-rate}_i + \varepsilon_{it}$

\vspace{0.15cm}
\scriptsize\color{inksoft} $\beta$ is identified off counties whose distance shrank more or less in the same transition period. The baseline-rate term corrects for mean reversion; SE clustered by state and, as a robustness check, by county.
\end{block}

\end{columns}
\end{frame}

\end{document}
